\section{引言与问题重述}
\subsection{问题背景}

我国南方山区地形起伏剧烈、沟谷深切，洪涝灾害发生后常出现道路塌方中断、电力与通信基站损毁的复合灾情。此时地面车辆难以在黄金救援期内抵达受灾村镇，“最后一公里”的物资投送与应急通信保障成为救援成败的关键。无人机具有起降场地要求低、受地形阻隔影响小、可快速部署的优点，已在应急物资投送、灾情侦察与临时通信中继等场景中得到应用\cite{yang2026dronehumanitarian,zhang2021truckdrone,wang2014reliefrouting}。

然而山区场景下的无人机运用面临\textbf{三重耦合困难}：\textbf{其一}，地形起伏直接决定航段必须爬升的高度，进而通过爬升能耗与载荷—航程关系影响单架次能力；\textbf{其二}，山体遮挡会造成视距链路中断，而应急通信要求运输无人机在爬升、巡航、下降与投送\textbf{全程}保持通信，必须引入中继无人机并合理选址；\textbf{其三}，机队规模、机型异构程度、共享电池数量与充电周转时间共同构成资源约束，使“能否按时送达”不再只是路径优化问题，而是\textbf{机队并行能力与调度规则共同决定}的资源—调度耦合问题。

本文以某山区乡镇的洪涝灾害救援为背景。临时调度中心 O01 同时是固定通信网关 G01 的所在地，下辖 15 个服务区，共需投送 80 个不可拆分的货箱，其中 30 个为“首批保障”货箱并附带更严格的截止时间。可用装备为 3 种型号共 8 架运输无人机、2 架中继无人机，以及按机型独立配置的共享电池与中继能源组件。区域 30\,m 数字高程模型（DEM）与通信链路预算参数由附件给出。

\subsection{问题重述}

\noindent \textbf{问题 1（单点往返运输能力与货箱组批）}：针对 3 种运输机型与 15 个服务区，确定每种机型在各服务区的\textbf{最大安全载荷}；在此基础上为每个服务区设计货箱组批方案，在满足载质量、装载体积与返航安全能量余量的前提下，综合往返架次数、总运输能耗与累计作业时间说明目标间的权衡关系，并分析返航安全余量 $\rhog$ 变化对结果的影响。

\noindent \textbf{问题 2（异构无人机多点多架次运输调度）}：允许单个架次访问一个或多个服务区，联合确定货箱组批、服务区访问顺序、运输机型、具体执行的实体无人机、共享电池分配与各架次开始时刻；在货箱不可拆分、载质量与装载体积、返航安全余量、无人机与共享电池数量、充电周转及物资时限等约束下，综合配送及时性、全部任务完成时间、运输能耗与架次数给出调度方案与指标权衡。

\noindent \textbf{问题 3（通信约束下的运输与中继联合调度）}：在问题 2 基础上引入\textbf{连续通信约束}——运输无人机在爬升、巡航、下降及物资投送阶段均应保持通信；与固定网关直连不可用时，可由空中中继无人机提供保障。需联合确定中继无人机的悬停位置、飞行海拔、服务时段与能源组件分配，并给出货箱交付、资源使用与通信保障情况。

\noindent \textbf{问题 4（救援任务分区与资源配置优化）}：以问题 3 的联合调度方案为基础，将 15 个服务区划分为 \textbf{2 组}与 \textbf{3 组}两种任务分区方案，保持问题 3 已确定的货箱组批、服务区访问顺序、运输与中继任务安排及通信保障关系不变；分别核算各任务组独立执行所需的运输无人机、共享电池、中继无人机与中继能源组件数量，并从配置规模、资源冗余、组间均衡与资源缺口四方面比较。

\subsection{本文的信息分级与证据状态}

应急优化类论文最易出现的失分点是\textbf{把“候选结果”写成“已证明的最优解”}。为此本文对全文出现的每一类数值明确标注其证据等级，并在正文中保持一致的表述：

\begin{enumerate}[label=\textbf{等级 \arabic*}]
	\item \textbf{附件核对的输入}：直接读取赛题附件得到的参数（机型参数、机队与电池数量、货箱清单、时限、DEM、链路预算参数）。这类数值可\textbf{直接引用}，且本文在 \S\ref{sec:data} 给出与附件汇总表的一致性核对。
	\item \textbf{本文复现并经验证的结果}：由本仓程序求解、并通过\textbf{独立校验器}复核的数值（问题一至问题四的架次、能耗、完工时间、资源占用等）。这类数值在正文中标注“本文结果”。
	\item \textbf{仅作对照的候选结果}：启发式得到的次优方案、或未附完整逐架次/逐箱明细的中间结果。这类数值\textbf{只用于权衡对比}，不声称最优性或资源可行性。
\end{enumerate}

相应地，本文对“最优性”的表述遵循\textbf{只在下界与可行构造同时成立时才称最优}的原则：问题一的 18 架次因逐服务区等于解析下界，故称“\textbf{架次数维度可证最优}”；问题二的 25 架次距理论下界 24 尚差 1 个，故只称“\textbf{经验证的可行解}”，不称最优；问题四的 $K=2/K=3$ 可行性来自精确图论计算，故称“\textbf{结构性结论}”。
